\begin{abstract}
SHAHID is a probabilistic artificial intelligence trust layer between smart-meter data and dispatch
decisions. When a group of meters under one distribution transformer stops reporting, it decides
whether the customers lost power or whether only their data stopped. A meter on the low-voltage side
of the transformer still sees the energy of customers whose data stopped, but not the energy of
customers who lost power. SHAHID compares that unexplained energy with the consumption it expects
from the silent customers, scores both causes with heavy-tailed likelihoods, and decides only when
the evidence is decisive; otherwise it abstains and hands the case to a person with the evidence, a
plain-language explanation and a tamper-evident log entry. All evidence is simulated on public German
benchmark networks. With thresholds frozen on one network and each evaluation run once, the method
called no labelled communication failure an outage on four other networks, missed at most one outage
in a hundred, and decided about half to three quarters of the events. On data with no events, random
single-meter data losses still produced about 15 to 33 false outage verdicts per transformer-year under the
registered configuration; a post-hoc gate reduced this to 0.2 to 1.0. The evaluations also measure
where the method fails: small silent groups, unmetered energy that follows the load, simultaneous
events and hourly readings. The method requires an energy meter at each transformer and does not
detect faults inside a single meter.
\end{abstract}

\begin{IEEEkeywords}
abstention, advanced metering infrastructure, communication failure, distribution networks, energy
balance, explainability, human oversight, outage detection, selective prediction, smart meters
\end{IEEEkeywords}
